\section{Formalization}
\label{sec: formalization}
In this section, we formalize key concepts in \ourmethod from a spatio-temporal perspective, and introduce the notations for clear clarification.
\subsection{Communication Topology}
In \ourmethod, a multi-agent system is modeled as a Spatio-Temporal Graph $\mathcal{G} = \{\mathcal{G}_t\}_{t=1}^T=\{(\mathcal{V}, \mathcal{E}_t,\mathrm{A}_t)\}_{t=1}^T$, where $t$ records the dialogue iteration. Note that the total iterations $T$ is flexible and not pre-defined. $\mathcal{V} = \{v_1, v_2, \cdots, v_N\}$ denotes $N$ agent nodes in MAS:
\begin{gather}
    v_i = \{\text{LLM}_i, \text{Profile}_i, \text{Tool}_i, \text{State}_i \},\label{eq1}
\end{gather}
where each agent $v_i$ is instantiated through a LLM with Profile (system prompts) and Tool, and its historical memory is maintained in State. During operation, each agent $v_t$ treats $\text{Profile}_i$ and $\text{State}_i$ as a prefix, and receives user prompts or responses from other agents in MAS.

In \ourmethod, the communication topology $\mathcal{G}_t$ changes over time, contributing to the varying communication link $\mathcal{E}_t$, which is described via the adjacent matrix $\mathrm{A}_t\in \{0,1\}^{N\times N}$, where $\mathrm{A}_t[i,j]=1$ denotes the edge $e_t[i,j] \in \mathcal{E}_t$ is connected and the information flows from $v_i$ to $v_j$, vice versa.

\subsection{Scheduling Pipeline}

Given a query $\mathcal{Q}$, \ourmethod first constructs the communication topology $\mathcal{G}_1$ from an initial anchor topology $\mathcal{G}_0$ for the first dialogue iteration. For simple questions, they can be easily tackled at once. For harder ones, multi-iteration dialogues are executed and fine-grained scheduling is planned through a Scheduler.  

In the $t$-th dialogue, the Scheduler considers the problem $\mathcal{Q}$, the iteration $t$, and the topology from the last iteration $\mathcal{G}_{t-1}$ to determine the topology of this iteration $\mathcal{G}_{t}$:
\begin{gather}
    \mathcal{G}_{t} = \text{Scheduler}(\mathcal{Q}, t, \mathcal{G}_{t-1})
\end{gather}
Then the communication topology $\mathcal{G}_t$ is compiled into a \textit{topological} execution order $\mathcal{S}_t$ of the $N$ agents in MAS:
\begin{gather}
    \mathcal{S}_t = [v^t_{s_1}, v^t_{s_2}, \cdots, v^t_{s_N}],\\
    \text{s.t.} \ \forall j>i, \ v^t_{s_j} \notin \mathcal{N}_{\text{in}}(v^t_{s_{i}})
\end{gather}
After \ourmethod executes $T$ iterations, the final solution is formed or partially formed inside the agent states, an aggregation operation like majority voting is applied to summarize for the output.